%% ma: L10-B14-0095
%% lop: 10
%% chuong: 5
%% bai: 14
%% dang: 10.14.7
%% loai: TLN
%% muc: VD
%% dap_an: "2,02"
%% nguon: "chuyen-de-cac-so-dac-trung-cua-mau-so-lieu-khong-ghep-nhom-toan-10-chuong-trinh-moi.pdf (GV. Phan Nhật Linh) (P05), Bài 03 Dạng 2 Phần III Câu 1"
%% kien_thuc: "Lớp 10 Bài 13 (số trung bình), Bài 14 (phương sai, độ lệch chuẩn)"
%% trang_thai: chinh-thuc
%% ghi_chu: "Đáp số gốc 2,025 chưa làm tròn đến hàng phần trăm theo yêu cầu đề; đúng là 2,02. Tự duyệt theo yêu cầu giáo viên 22/10/2026."
\begin{ex}
Số liệu sau đây cho biết số con được sinh ra trong $20$ hộ gia đình được khảo sát ở một địa phương:
\begin{center}
\begin{tabular}{|c|c|c|c|c|c|c|c|c|c|}\hline
$2$ & $2$ & $3$ & $5$ & $2$ & $4$ & $3$ & $2$ & $1$ & $9$\\\hline
$5$ & $3$ & $2$ & $4$ & $1$ & $0$ & $3$ & $2$ & $1$ & $6$\\\hline
\end{tabular}
\end{center}
Tìm độ lệch chuẩn của mẫu số liệu (làm tròn đến hàng phần trăm).
\shortans{2,02}
\loigiai{$\overline x=\dfrac{60}{20}=3$; $s^2=\dfrac{1}{20}\sum\left(x_i-3\right)^2=4{,}1$; $s=\sqrt{4{,}1}\approx2{,}02$.}
\end{ex}
